% The picture gallery: every picture of pythonfigs.sty with the program it
% depicts, in the layout of the decks' notes (the same class options and
% preamble as notes.tex), for the author's review before any deck uses
% the pictures.
\documentclass[a4paper,10pt,oneside,openright,
               oldfontcommands]{memoir}
../../containers/slides-dicts/preamble.tex

\usepackage[noamsthm,notheorems]{beamerarticle}
\usepackage{pythonfigs}

% The gallery's pictures, one macro each, so that gallery.tex (the notes
% layout) and gallery-slides.tex (beamer) draw exactly the same picture.
% The running example is a bank, as in the recap.

% balance = 1000; old = balance; balance = balance + 500: three moments,
% boxes only.
\newcommand{\figBoxSteps}{%
  \begin{pyfig}
    \pynames[anchor=north west, at={(0,0)}]{s1}{balance: 1000}
    \pynames[anchor=north west, at={(32mm,0)}]{s2}{balance: 1000, old: 1000}
    \pynames[anchor=north west, at={(64mm,0)}]{s3}{balance: 1500, old: 1000}
    \foreach \s/\t in {s1/1, s2/2, s3/3}
      \node[pyfig region title] at ([yshift=2mm]\s.north west)
        {Efter rad \t};
  \end{pyfig}%
}

% owner = "Malvina"
\newcommand{\figString}{%
  \begin{pyfig}
    \pyvar{owner}{"Malvina"}
  \end{pyfig}%
}

% transactions = [500, -200, 1000]; history = transactions;
% backup = transactions.copy()
\newcommand{\figAlias}{%
  \begin{pyfig}
    \pynames[row sep=5mm]{n}%
      {transactions: \pyptr, history: \pyptr, backup: \pyptr}
    \pylist[right=12mm of n-transactions]{l1}{500, -200, 1000}
    \pylist[right=12mm of n-backup]{l2}{500, -200, 1000}
    \pyarrow{n-transactions}{l1}
    \pyarrow{n-history}{l1.west}
    \pyarrow{n-backup}{l2}
  \end{pyfig}%
}

% A list of strings with its indices (a construct in general).
\newcommand{\figListStrings}{%
  \begin{pyfig}
    \pyvar{names}{\pyptr}
    \pylist[right=12mm of names]{l}{"Ronja", "Pippi", "Madicken"}
    \pyarrow{names}{l}
  \end{pyfig}%
}

% A tuple: one transaction, (date, amount).
\newcommand{\figTuple}{%
  \begin{pyfig}
    \pyvar{transaction}{\pyptr}
    \pytuple[right=12mm of transaction]{t}{"2026-10-26", 500}
    \pyarrow{transaction}{t}
  \end{pyfig}%
}

% A nested list: the outer list's cells refer to the inner lists.
\newcommand{\figNested}{%
  \begin{pyfig}
    \pyvar{grid}{\pyptr}
    \pylist[right=12mm of grid]{o}{\pyptr, \pyptr}
    \pylist[below=10mm of o-0.south west, anchor=north east,
            xshift=2mm, pyfig type at=left]{r0}{1, 2}
    \pylist[below=10mm of o-1.south east, anchor=north west,
            xshift=-2mm, pyfig type at=left]{r1}{3, 4}
    \pyarrow{grid}{o}
    \pyarrow{o-0}{r0}
    \pyarrow{o-1}{r1}
  \end{pyfig}%
}

% A list of (date, amount) pairs.
\newcommand{\figPairs}{%
  \begin{pyfig}
    \pyvar{transactions}{\pyptr}
    \pylist[right=12mm of transactions]{o}{\pyptr, \pyptr}
    \pytuple[below=9mm of o-0.south west, anchor=north west,
             pyfig type at=left]{t0}%
      {"2026-10-26", 500}
    \pytuple[below=7mm of t0.south west, anchor=north west,
             pyfig type at=left]{t1}%
      {"2026-10-27", -200}
    \pyarrow{transactions}{o}
    \pyarrow{o-0}{t0.north -| o-0}
    \pyarrow[out=0, in=0, looseness=1.6]{o-1}{t1.east}
  \end{pyfig}%
}

% A dict: owner -> balance.
\newcommand{\figDict}{%
  \begin{pyfig}
    \pyvar{accounts}{\pyptr}
    \pydict[right=12mm of accounts]{d}{"ronja": 1500, "pippi": 300}
    \pyarrow{accounts}{d}
  \end{pyfig}%
}

% A dict whose values are dicts: account number -> the account.
\newcommand{\figDictOfDicts}{%
  \begin{pyfig}
    \pyvar{accounts}{\pyptr}
    \pydict[right=10mm of accounts]{d}%
      {"1234": \pyptr, "5678": \pyptr, "9012": \pyptr}
    \pydict[right=12mm of d-0, yshift=9mm]{r}%
      {"owner": "Ronja", "balance": 1500}
    \pydict[right=12mm of d-1]{p}%
      {"owner": "Pippi", "balance": 300}
    \pydict[right=12mm of d-2, yshift=-9mm]{m}%
      {"owner": "Madicken", "balance": 2000}
    \pyarrow{accounts}{d}
    \pyarrow{d-0}{r}
    \pyarrow{d-1}{p}
    \pyarrow{d-2}{m}
  \end{pyfig}%
}

% A set of owners.
\newcommand{\figSet}{%
  \begin{pyfig}
    \pyvar{owners}{\pyptr}
    \pyset[right=12mm of owners]{s}{"Ronja", "Pippi", "Madicken"}
    \pyarrow{owners}{s}
  \end{pyfig}%
}

% A list used as a stack: the transactions to undo.
\newcommand{\figStack}{%
  \begin{pyfig}
    \pyvar{undo}{\pyptr}
    \pystack[right=12mm of undo]{s}{500, -200, 1000}
    \pyarrow{undo}{s}
  \end{pyfig}%
}

% A deque used as a queue: the customers waiting.
\newcommand{\figQueue}{%
  \begin{pyfig}
    \pyqueue{q}{"Ronja", "Pippi", "Madicken"}
    \pyvar[above=7mm of q-1]{customers}{\pyptr}
    \pyarrow{customers}{q-1.north}
  \end{pyfig}%
}

% An Account object and two names for it.
\newcommand{\figObject}{%
  \begin{pyfig}
    \pynames[row sep=5mm]{n}{account: \pyptr, mine: \pyptr}
    \pyobject[right=15mm of n-account, yshift=-4mm]{a}{Account}%
      {owner: "Ronja", number: "1234", balance: 1500}
    \pyarrow{n-account}{a}
    \pyarrow{n-mine}{a}
  \end{pyfig}%
}

% A Bank object: object -> dict -> objects.
\newcommand{\figBank}{%
  \begin{pyfig}
    \pyvar{bank}{\pyptr}
    \pyobject[right=10mm of bank]{b}{Bank}%
      {name: "Malvinas bank", accounts: \pyptr}
    \pydict[below=10mm of b.south east, anchor=north east]{d}%
      {"1234": \pyptr, "5678": \pyptr}
    \pyobject[right=9mm of d-0, yshift=8mm]{a0}{Account}%
      {owner: "Ronja", number: "1234", balance: 1500}
    \pyobject[right=9mm of d-1, yshift=-14mm]{a1}{Account}%
      {owner: "Pippi", number: "5678", balance: 300}
    \pyarrow{bank}{b}
    \pyarrow[out=-90, in=90]{b-accounts}{d}
    \pyarrow{d-0}{a0}
    \pyarrow{d-1}{a1}
  \end{pyfig}%
}

% A call: main has called deposit(account, 500).
\newcommand{\figCall}{%
  \begin{pyfig}
    \pyframe{m}{main}{account: \pyptr}
    \pyframe[below=9mm of m.south west, anchor=north west]{dep}%
      {deposit}{account: \pyptr, amount: 500}
    \pydict[right=10mm of m-account, yshift=-8mm]{d}%
      {"owner": "Ronja", "balance": 1500}
    \pyarrow{m-account}{d}
    \pyarrow{dep-account}{d}
  \end{pyfig}%
}

% Scope: a global balance and a local balance.
\newcommand{\figScope}{%
  \begin{pyfig}
    \pyframe*{g}{Globala namn}{balance: 1000}
    \pyframe[below=9mm of g.south west, anchor=north west]{f}%
      {add_interest}{rate: 0.02, balance: 1020.0}
  \end{pyfig}%
}

% if/elif/else: a withdrawal.
\newcommand{\figIf}{%
  \pyflowif{amount <= 0: print("Ogiltigt belopp"),
            amount > balance: print("Otillräckligt saldo")}%
    {balance -= amount}%
}

% while: asking for the PIN until it is right.
\newcommand{\figWhile}{%
  \pyflowwhile{pin != "1234"}{pin = input("PIN: ")}%
}

% Saving the accounts to a file, and reading them back.
\newcommand{\figFile}{%
  \begin{pyfig}
    \pynames[row sep=10mm]{n}{accounts: \pyptr, rader: \pyptr}
    \pydict[right=10mm of n-accounts]{d}{"ronja": 1500, "pippi": 300}
    \pylist[right=10mm of n-rader]{l}{{"ronja,1500\n"}, {"pippi,300\n"}}
    \pyarrow{n-accounts}{d}
    \pyarrow{n-rader}{l}
    \pyfile[below right=14mm and 8mm of l.south east, anchor=north west]{f}%
      {konton.txt}%
      {ronja,1500 \\ pippi,300}
    \pyio[out=0, in=90]{d.east}{f.north east}{write}
    \pyio{f.north west}{l.south east}{readlines}
    \pyregion{(n)(d)(l)}{Programmet (minnet)}
    \pyregion{(f)}{Disken}
  \end{pyfig}%
}


% \entry{<label>}{<caption>}{<picture macro>}: the picture as a figure
% with a side caption, as a deck's notes would set it.  The program it
% depicts stands before it, in a minted environment.
% Each entry starts where at least two fifths of the page are left, so
% that its program and its picture stay on one page.
\newcommand{\entrysection}[1]{\needspace{0.4\textheight}\section{#1}}
\newcommand{\entry}[3]{%
  \begin{figure}[htbp]
    \begin{sidecaption}{#2}[fig:#1]
      #3
    \end{sidecaption}
  \end{figure}
}

\begin{document}
\title{Bildspråket: ett galleri}
\author{Daniel Bosk}
\date{Utkast v1, 10 oktober 2026}
\maketitle

\begin{abstract}
  Varje bild i paketet \texttt{pythonfigs}, med programmet den visar och
  en bildtext.  Modellen är beslutet från den 10 oktober (alternativ b):
  ett namn som håller ett tal eller en sträng är en låda med värdet i;
  från listorna och framåt ritas objektet för sig och namnets låda håller
  en pil till det.  I en behållare står tal och strängar i cellen, och en
  behållare i en behållare är en pil.  Exemplen kommer ur repetitionens
  genomgående exempel, banken; där en bild visar en konstruktion i
  allmänhet är exemplet fristående.  Sist samma bilder på riktiga
  bilder (\texttt{beamer}), så att storleken syns.
\end{abstract}

\clearpage
\tableofcontents*

\chapter{Galleriet}

\entrysection{Namn och värden}

Tre rader, och lådorna efter varje rad:
\begin{minted}{python}
balance = 1000
old = balance
balance = balance + 500
\end{minted}
\Cref{fig:steps} visar att \mintinline{python}{old = balance} lägger en
kopia av värdet i den nya lådan: när \mintinline{python}{balance} sedan
får ett nytt värde ligger 1000 kvar i \mintinline{python}{old}.
\entry{steps}{Tilldelning med lådor: namnen med sina värden efter
  varje rad.}{\figBoxSteps}

\entrysection{En sträng}

\begin{minted}{python}
owner = "Malvina"
\end{minted}
\Cref{fig:string} ritar strängen som ett tal: värdet står i lådan.
\entry{string}{En sträng i en låda, citattecknen med.}{\figString}

\entrysection{Två namn på en lista, och en kopia}

\begin{minted}{python}
transactions = [500, -200, 1000]
history = transactions
backup = transactions.copy()
\end{minted}
I \cref{fig:alias} pekar \mintinline{python}{transactions} och
\mintinline{python}{history} på samma lista, medan
\mintinline{python}{backup} pekar på en lista av sin egen.
\entry{alias}{En lista och två namn (\texttt{history = transactions})
  mot två listor (\texttt{backup = transactions.copy()}).}{\figAlias}

\entrysection{En lista av strängar}

\begin{minted}{python}
names = ["Ronja", "Pippi", "Madicken"]
\end{minted}
\Cref{fig:liststrings} visar cellerna med sina index under.
\entry{liststrings}{En lista av strängar; under varje cell dess
  index.}{\figListStrings}

\entrysection{En tupel}

\begin{minted}{python}
transaction = ("2026-10-26", 500)
\end{minted}
En tupel (\cref{fig:tuple}) ritas som en lista men med runda hörn och
typen \texttt{tuple}.
\entry{tuple}{En tupel, en transaktion som datum och belopp: rundade
  hörn som de runda parenteserna, typen ovanför.}{\figTuple}

\entrysection{En lista av listor}

\begin{minted}{python}
grid = [[1, 2], [3, 4]]
\end{minted}
De inre listorna är objekt för sig, och cellerna i den yttre pekar på
dem (\cref{fig:nested}).
\entry{nested}{En lista av listor: den yttre listans celler pekar på
  de inre.}{\figNested}

\entrysection{En lista av par}

\begin{minted}{python}
transactions = [("2026-10-26", 500), ("2026-10-27", -200)]
\end{minted}
\Cref{fig:pairs} är samma slags bild som \cref{fig:nested}, med tupler
i stället för listor.
\entry{pairs}{En lista av par (datum, belopp): varje cell pekar på en
  tupel.}{\figPairs}

\entrysection{En uppslagslista}

\begin{minted}{python}
accounts = {"ronja": 1500, "pippi": 300}
\end{minted}
\Cref{fig:dict} ritar uppslagslistan som en tabell, nyckel till vänster
och värde till höger.
\entry{dict}{En uppslagslista från ägare till saldo: en rad per nyckel,
  nyckeln till vänster, värdet till höger.}{\figDict}

\entrysection{En uppslagslista av uppslagslistor}

\begin{minted}{python}
accounts = {
    "1234": {"owner": "Ronja", "balance": 1500},
    "5678": {"owner": "Pippi", "balance": 300},
    "9012": {"owner": "Madicken", "balance": 2000},
}
\end{minted}
I \cref{fig:dictofdicts} är varje värde i den yttre uppslagslistan en
pil till en inre.
\entry{dictofdicts}{Uppslagslistor i en uppslagslista: kontonumret
  pekar på kontots uppgifter.}{\figDictOfDicts}

\entrysection{En mängd}

\begin{minted}{python}
owners = {"Ronja", "Pippi", "Madicken"}
\end{minted}
En mängd har ingen ordning och inga index, så elementen i
\cref{fig:set} står huller om buller.
\entry{set}{En mängd av ägare: elementen utan ordning och utan
  index.}{\figSet}

\entrysection{En stack}

\begin{minted}{python}
undo = []
undo.append(500)
undo.append(-200)
undo.append(1000)
\end{minted}
\Cref{fig:stack} är listan med båda operationerna i samma ände: den
transaktion som ångras först är den som gjordes sist.
\entry{stack}{En lista som stack: \texttt{append} lägger till och
  \texttt{pop} tar bort i samma ände.}{\figStack}

\entrysection{En kö}

\begin{minted}{python}
customers = collections.deque()
customers.append("Ronja")
customers.append("Pippi")
customers.append("Madicken")
\end{minted}
\Cref{fig:queue} har operationerna i var sin ände: den kund som kom
först blir betjänad först.
\entry{queue}{En \texttt{deque} som kö: \texttt{append} lägger till i
  slutet, \texttt{popleft} tar bort först.}{\figQueue}

\entrysection{Ett objekt och två namn}

\begin{minted}{python}
account = Account("Ronja", "1234", 1500)
mine = account
\end{minted}
Ett objekt (\cref{fig:object}) ritas med klassen som typ och ett namn
per attribut, som namnen i en funktion.
\entry{object}{Ett objekt med tre attribut, och två namn som pekar på
  det.}{\figObject}

\entrysection{Objekt i objekt}

\begin{minted}{python}
bank = Bank("Malvinas bank")
bank.add(Account("Ronja", "1234", 1500))
bank.add(Account("Pippi", "5678", 300))
\end{minted}
Banken håller sina konton i en uppslagslista från kontonummer till
konto: från objektet via uppslagslistan till objekten
(\cref{fig:bank}).
\entry{bank}{En bank med två konton: attributet \texttt{accounts} pekar
  på en uppslagslista, och den pekar på kontona.}{\figBank}

\entrysection{Ett anrop}

\begin{minted}{python}
def deposit(account, amount):
    account["balance"] = account["balance"] + amount


def main():
    account = {"owner": "Ronja", "balance": 1500}
    deposit(account, 500)


main()
\end{minted}
\Cref{fig:call} visar läget inne i \mintinline{python}{deposit}, innan
saldot ändras: en ruta per anrop som pågår, den senaste nederst.
Parametern \mintinline{python}{account} pekar på samma uppslagslista
som argumentet, medan \mintinline{python}{amount} har sin egen låda.
\entry{call}{Två anrop som pågår: \texttt{main} har anropat
  \texttt{deposit}, och båda namnen \texttt{account} pekar på samma
  uppslagslista.}{\figCall}

\entrysection{Räckvidd}

\begin{minted}{python}
balance = 1000


def add_interest(rate):
    balance = 1000 * (1 + rate)
    print(balance)


add_interest(0.02)
print(balance)
\end{minted}
Inne i \mintinline{python}{add_interest} finns två
\mintinline{python}{balance}, ett globalt och ett lokalt
(\cref{fig:scope}).  Bilden visar också vad
\mintinline{python}{def} gör: \mintinline{python}{add_interest} är ett
globalt namn som pekar på ett funktionsobjekt.  Det är den enda bilden
som ritar funktionerna; i de andra utelämnas de, för att bilderna ska
hållas små.
\entry{scope}{Ett globalt och ett lokalt namn med samma namn, i var
  sin ruta.  Funktionen är själv ett namn, som pekar på ett objekt av
  typen \texttt{function}.}{\figScope}

\entrysection{Ett returvärde}

\begin{minted}{python}
def interest(balance, rate):
    return balance * rate


def main():
    due = interest(1500, 0.02)
    print(due)


main()
\end{minted}
\Cref{fig:return} visar anropet i två ögonblick: strax innan
\mintinline{python}{interest} returnerar, med returvärdet i sin ram, och
efter anropet, när ramen är borta och värdet har hamnat i
\mintinline{python}{due}.
\entry{return}{Ett returvärde: först i den anropade funktionens ram,
  sedan i namnet hos den som anropade.}{\figReturn}

\entrysection{Villkor}

\begin{minted}{python}
if amount <= 0:
    print("Ogiltigt belopp")
elif amount > balance:
    print("Otillräckligt saldo")
else:
    balance -= amount
\end{minted}
\Cref{fig:if} följer programmet uppifrån: det första villkoret som är
sant avgör vägen.
\entry{if}{\texttt{if}, \texttt{elif} och \texttt{else} som
  flödesschema för ett uttag: villkoren nedåt, en gren åt höger för
  varje sant villkor.}{\figIf}

\entrysection{Upprepning}

\begin{minted}{python}
pin = input("PIN: ")
while pin != "1234":
    pin = input("PIN: ")
\end{minted}
I \cref{fig:while} går vägen tillbaka till villkoret efter varje varv.
\entry{while}{\texttt{while} som flödesschema: kroppen, och vägen
  tillbaka till villkoret.}{\figWhile}

\entrysection{Upprepning med \texttt{for}}

\begin{minted}{python}
transactions = [500, -200, 1000]
balance = 1000
for amount in transactions:
    balance = balance + amount
\end{minted}
\Cref{fig:for} ritar slingan som flödesschema, med varven i en tabell
bredvid: varje varv tar nästa värde ur listan, och slingan slutar när
värdena tar slut.
\entry{for}{\texttt{for} som flödesschema, och samma slinga varv för
  varv: \texttt{amount} och \texttt{balance} efter varje
  varv.}{\figFor\hspace{1.5em}\tabFor}

\entrysection{Spara och läsa en fil}

\begin{minted}{python}
with open("konton.txt", "w") as fil:
    for owner, balance in accounts.items():
        fil.write(f"{owner},{balance}\n")

with open("konton.txt") as fil:
    rader = fil.readlines()
\end{minted}
\Cref{fig:file} skiljer programmets värden i minnet från filen på
disken.
\entry{file}{Programmet och filen: \texttt{write} skriver kontona till
  filen, \texttt{readlines} hämtar raderna till en lista.}{\figFile}

\entrysection{Egentligen är alla namn pilar}

Bilderna ovan ritar tal och strängar som lådor.  I Python är också de
objekt, och alla namn är pilar.  Det här avsnittet ritar om några av
bilderna så: ett tal är ett objekt med sin typ ovanför, som en lista.
Kursen ritar ändå lådor först, eftersom ett tal eller en sträng inte
kan ändras: då ger lådan och pilen alltid samma resultat, och lådan är
enklare att följa tills det första objektet som kan ändras, listan,
gör skillnaden synlig.

\subsection{Tilldelning}

\begin{minted}{python}
balance = 1000
old = balance
balance = balance + 500
\end{minted}
I \cref{fig:arrowsteps} pekar båda namnen först på samma tal.  Den
tredje raden skapar ett nytt tal och flyttar \mintinline{python}{balance}
dit, medan \mintinline{python}{old} pekar kvar på 1000.  Resultatet är
detsamma som i \cref{fig:steps}, av ett annat skäl.
\entry{arrowsteps}{Tilldelning med pilar: efter varje rad.  Ingen
  kopia görs; den tredje raden flyttar en pil.}{\figArrowSteps}

\needspace{0.4\textheight}
\subsection{Ett anrop}

Samma anrop som i \cref{fig:call}, \mintinline{python}{deposit(account,
500)}, med pilar (\cref{fig:arrowcall}): båda parametrarna lämnas över
på samma sätt, som en pil till objektet hos den som anropar.  Också
talet i uppslagslistan är ett objekt: cellen för
\mintinline{python}{"balance"} håller en pil till det.
\entry{arrowcall}{Ett anrop med pilar: \texttt{account} och
  \texttt{amount} pekar båda på anroparens objekt.  Saldot är ett objekt
  för sig, som uppslagslistan pekar på.}{\figArrowCall}

\needspace{0.4\textheight}
\subsection{Att flytta en pil och att ändra ett objekt}

Varför ser ett tal ut att kopieras när en lista inte gör det?
\begin{minted}{python}
def deposit(account, amount):
    amount = amount + 1
    account["balance"] += amount
    account = {"owner": "Pippi", "balance": 0}


def main():
    amount = 500
    account = {"owner": "Ronja", "balance": 1500}
    deposit(account, amount)
    print(account["balance"], amount)


main()
\end{minted}
Programmet skriver ut \texttt{2001 500}.  Den första raden i
\mintinline{python}{deposit} flyttar den lokala pilen
\mintinline{python}{amount} till ett nytt tal, och anroparen märker
ingenting (\cref{fig:rebindA}).  Den andra raden ändrar objektet som
båda \mintinline{python}{account} pekar på: pilen i dess cell för
\mintinline{python}{"balance"} flyttas från talet 1500 till ett nytt
tal, 2001.  Båda namnen pekar kvar på samma uppslagslista, så det
märker anroparen (\cref{fig:rebindB}).  Den tredje raden flyttar den lokala pilen
\mintinline{python}{account} till en ny uppslagslista, och det märker
anroparen inte heller (\cref{fig:rebindC}).  Gränsen går alltså mellan
att flytta en pil och att ändra ett objekt, inte mellan typer: tal och
strängar ser bara ut att kopieras, eftersom de inte kan ändras.
\entry{rebindA}{Efter \texttt{amount = amount + 1}: den lokala pilen
  pekar på ett nytt tal, \texttt{main} pekar kvar på 500.}{\figRebindA}
\entry{rebindB}{Efter \mbox{\texttt{account[\textquotedbl
  balance\textquotedbl] += amount}}: pilen i den gemensamma
  uppslagslistan pekar nu på 2001, och \texttt{main} ser det.  Talet 1500
  har ingen pil kvar.}{\figRebindB}
\entry{rebindC}{Efter \texttt{account = \{...\}}: den lokala pilen pekar
  på en ny uppslagslista, \texttt{main} pekar kvar på den
  gamla.}{\figRebindC}

\chapter{Samma bild på en bild}

Fyra av bilderna på riktiga bilder i föreläsningarnas
\texttt{beamer}-tema (\cref{fig:slide1,fig:slide2,fig:slide3,fig:slide4}).

\begin{figure}[htbp]
  \centering
  \includegraphics[page=1,width=0.9\textwidth]{gallery-slides.pdf}
  \caption{Två namn på en lista och en kopia, programmet ovanför.}
  \label{fig:slide1}
\end{figure}

\begin{figure}[htbp]
  \centering
  \includegraphics[page=2,width=0.9\textwidth]{gallery-slides.pdf}
  \caption{Kontona som uppslagslista av uppslagslistor, programmet
    ovanför.}
  \label{fig:slide2}
\end{figure}

\begin{figure}[htbp]
  \centering
  \includegraphics[page=3,width=0.9\textwidth]{gallery-slides.pdf}
  \caption{Ett anrop: anropsramarna bredvid en förklaring.}
  \label{fig:slide3}
\end{figure}

\begin{figure}[htbp]
  \centering
  \includegraphics[page=4,width=0.9\textwidth]{gallery-slides.pdf}
  \caption{\texttt{for} som flödesschema, med programmet och varven
    bredvid.}
  \label{fig:slide4}
\end{figure}

\end{document}
